\section{Introduction}
\label{sec:intro}

Fine-tuning is the standard route for injecting new knowledge into pretrained
language models, whether through supervised instruction tuning, domain
adaptation, or parameter-efficient updates \citep{hu2022lora}. A growing body
of evidence shows that this process memorizes far more readily than it
generalizes: models fit training-format facts quickly
\citep{tirumala2022memorization,carlini2023quantifying}, while behavioral
control of the acquired knowledge---answering paraphrases, chaining facts,
or refusing when unsure---arrives late or not at all
\citep{kandpal2023large,allen2023physics}. For practitioners this gap is
expensive: a model that \emph{knows} an entity (its answer is recoverable
from its internal states) but cannot \emph{say} it under mild distribution
shift will silently fail downstream, and standard observables cannot
distinguish the two regimes.

We propose to make this distinction measurable \emph{during} training, and
to act on it. Our approach, MLPR (Mid-Layer Probe Regularization), wraps a
parameter-efficient fine-tuning run with three coupled instruments. First,
a small linear probe $p_\phi$ attached to a mid layer $l^\star$ reads
entity identities from the backbone's hidden states; because the probe
trains on detached states during a warm-up phase, it is a pure observer
that cannot steer the backbone. Second, a teacher-forced likelihood metric
estimates recall saturation $\amemlf$ each epoch by scoring the answer
tokens given the training-format question---a ``knowing'' signal that, unlike
free-generation exact match, does not conflate knowledge with answer
extraction. Third, a saturation gate compares a monitored signal against a
threshold $\tau_0$ and, once crossed, ramps in a probe-alignment loss with
weight $\lambda(t)$ so that the backbone is regularized toward keeping acquired
identities linearly decodable at the mid layer. An eval-aware adaptive
controller completes the loop by damping the learning rate when held-out
loss degrades, subject to an explicit floor that keeps memorization
dynamics alive.

We instantiate MLPR on Qwen2.5-0.5B-Instruct and Qwen3-1.7B
\citep{qwen2024qwen25,qwen3}, training a grid of ten fully
instrumented runs (24 epochs; two 16-epoch censored replicates) on a
synthetic entity-question corpus with
a 1{,}600-entity vocabulary (per-step losses, per-epoch gate
evaluations, per-sample final evaluations, adapters and probe
checkpoints). The grid crosses two gate signals---output-likelihood
recall vs.\ representation-space probe decodability---with three
engagement schedules (sustained, decayed-to-zero, output-triggered),
completing the engagement-schedule arm that was budget-censored in our
earlier large-scale study. The resulting dataset of training dynamics
yields findings we summarize here and develop in
Section~\ref{sec:results}:

\begin{itemize}[leftmargin=1.4em,itemsep=1pt,topsep=2pt]
  \item \textbf{The gate signal decides when engagement is possible}
  (Section~\ref{sec:res:sat}). The probe's representation decodability
  crosses its threshold at epoch 3 (0.5B) and epochs 4--6 (1.7B), while
  the output-likelihood signal crosses $\tau_0{=}0.8$ only at epoch
  17--19---and in one output-gated replicate never crosses at all.
  Gating on the representation is scale-aware by construction; gating
  on the output is reachable at scale but unstable in crossing time.
  \item \textbf{The engagement schedule decides whether memorization
  completes} (Section~\ref{sec:res:sched}). Under identical data, seed,
  and optimizer, final recall spans $0.012$--$0.934$ at 1.7B. Late
  output-triggered engagement reaches SFT parity ($0.926$ vs.\ $0.934$,
  McNemar $p{=}0.83$) and completes the regularized phase---8 stable
  epochs at full probe weight---that our earlier large-scale run left
  budget-censored. Early sustained injection is high-variance
  ($0.109$, $0.805$); transient injection that decays to zero
  collapses recall to $0.012$ while the probe still decodes $94.8\%$
  of entities.
  \item \textbf{Knowing and saying dissociate at both scales}
  (Section~\ref{sec:res:gap}). Teacher-forced recall reaches
  $0.91$--$0.93$ and token-level recall $0.99$, yet free-generation EM
  stays $\le 0.012$ and generalization EM $\le 0.004$ in \emph{every}
  arm: the saying gap is invariant across a $3.4\times$ parameter
  range and across all schedules.
  \item \textbf{Probes detect memorization early and cheaply}
  (Section~\ref{sec:res:control}). The probe crosses $0.5$ decodability
  between epochs 4 and 7 everywhere, leading output recall by 8--18
  epochs (and leading arms whose output never crosses at all), with
  loss two orders of magnitude below chance.
  \item \textbf{Updates concentrate in upper-layer MLPs}
  (Section~\ref{sec:res:anatomy}). $\Delta W = \tfrac{\alpha}{r}BA$ per
  LoRA-adapted projection puts $65$--$71\%$ of update mass in MLP
  projections, key/value attention nearly untouched, consistent with
  factual storage accounts \citep{meng2022rome,geva2021ffn}.
\end{itemize}

Our contributions are: (i) a training-time instrumentation stack---probe,
likelihood-based recall metric, saturation gate with selectable signal,
and engagement schedules---that turns the knowing--using gap into an
observable, thresholdable quantity; (ii) a completed
engagement-schedule grid at two scales with paired McNemar analysis,
providing the first fully exercised regularized phase (8 epochs at
full probe weight) and its outcome parity with SFT; (iii) a
dissociation analysis separating representation decodability,
teacher-forced recall, free-generation EM, and chaining/intersection
transfer, including a representation-without-output regime induced by
transient engagement; and (iv) engineering lessons from instrumenting
real runs, including a documented failure mode and repair for
eval-aware learning-rate controllers (Appendix~\ref{app:spiral}).
All telemetry, code, and trained adapters are released.

\section{Related Work}
\label{sec:related}

\textbf{Probing internal representations.} Linear probes are a standard
instrument for asking what models encode
\citep{alain2017probes,belinkov2022probing}; control tasks guard against
probes that exploit surface statistics \citep{hewitt2019control}. Our
probe differs in role: trained \emph{alongside} the run on detached
states, its readings drive a control decision---activating a
regularization pathway---turning it from a diagnostic into a controller
input. Richer readouts \citep{bricken2023monosemanticity} could
substitute in future work.

\textbf{Knowledge storage and memorization dynamics.} Mechanistic and
editing studies localize factual knowledge in mid-to-upper MLP blocks
that behave as key--value memories
\citep{meng2022rome,meng2023memit,geva2021ffn,dai2022knowledgeneurons,
mitchell2022mend}; our anatomy analysis is consistent with this picture
from the direction of training dynamics. Synthetic entity corpora with
known ground truth have been used to study capacity and retrieval
\citep{allen2023physics}; our generator follows this tradition but
exposes per-epoch recall dynamics. Memorization of training data is
well documented \citep{carlini2023quantifying,zhang2017understanding},
large models memorize more before they overfit in loss terms
\citep{tirumala2022memorization}, and knowledge acquisition is strongly
scale-dependent \citep{kaplan2020scaling,kandpal2023large}; our grid
gives a fine-grained per-epoch view in the fine-tuning regime. The
dissociation we measure extends paraphrase-consistency findings
\citep{elazar2021measuring,mallen2023pararel} into training dynamics,
and recalls the competence--expression decoupling of grokking
\citep{power2022grokking}, here driven by data format rather than
weight decay.

\textbf{Adaptive training control.} Adaptive learning-rate schedules
respond to loss plateaus or gradient statistics; our controller adds a
held-out-loss criterion with a twist specific to the memorization
regime: because rising held-out loss is the very signal the gate waits
for, damping must be bounded lest the controller extinguish the
phenomenon it monitors. We are not aware of prior work analyzing this
controller--phenomenon interaction; we document a negative result (an
unbounded decay spiral) and its repair.
